\documentclass[reqno]{amsart}
\usepackage[T1]{fontenc}
\usepackage{amsmath,amssymb,amsthm}
\usepackage{xr-hyper}
\usepackage[hidelinks]{hyperref}
\externaldocument[ch-]{20261125T003000--the-chain-and-the-square-root__workbook}
\externaldocument[sw-]{20261125T015500--siegel-weil-in-the-smallest-case__workbook}
\externaldocument[tn-]{20261125T002100--definite-theta-lifts-and-norms__workbook}
% Chapter number; \ref{chap} gives it.
\newcounter{chap}
% Chapter numbers, generated from 20261125T012000--grh-workbook__guide.org
% by scripts/make-chapters.py.  Do not edit: change the order of the
% index and rerun the script (./preview-site.sh and CI do this).
%
% Usage in a sheet, after \newcounter{chap}:
%   % Chapter numbers, generated from 20261125T012000--grh-workbook__guide.org
% by scripts/make-chapters.py.  Do not edit: change the order of the
% index and rerun the script (./preview-site.sh and CI do this).
%
% Usage in a sheet, after \newcounter{chap}:
%   % Chapter numbers, generated from 20261125T012000--grh-workbook__guide.org
% by scripts/make-chapters.py.  Do not edit: change the order of the
% index and rerun the script (./preview-site.sh and CI do this).
%
% Usage in a sheet, after \newcounter{chap}:
%   \input{chapters}\chapnumber{<sheet ID>}
\makeatletter
\providecommand\chapters@set[2]{\@namedef{chapters@#1}{#2}}
\chapters@set{20261125T001000}{1}% zeros-to-ramp-sums
\chapters@set{20261125T001100}{2}% inert-primes-and-estimate-b
\chapters@set{20261125T005000}{3}% galois-fourier-and-averaging
\chapters@set{20261125T005500}{4}% columns-from-a-sphere
\chapters@set{20261125T000000}{5}% slope-method
\chapters@set{20261126T030100}{6}% a-twist-for-each-prime
\chapters@set{20261125T003000}{7}% the-chain-and-the-square-root
\chapters@set{20261125T010000}{8}% rank-lemma-and-zero-estimate
\chapters@set{20261125T010500}{9}% poincare-series-and-near-orthogonality
\chapters@set{20261124T223000}{10}% elliptic-curves-for-the-proof
\chapters@set{20261125T009500}{11}% fock-model-and-theta-intertwiners
\chapters@set{20261124T223100}{12}% complex-multiplication-and-inert-reduction
\chapters@set{20261125T008000}{13}% one-cm-determinant
\chapters@set{20261124T214500}{14}% supersingular-theta-divisibility
\chapters@set{20261124T214600}{15}% the-isotropic-plane
\chapters@set{20261125T013700}{16}% pairings-on-a-formal-group-of-height-two
\chapters@set{20261125T013500}{17}% hecke-nondegeneracy-on-the-tree
\chapters@set{20261126T030200}{18}% the-bound-is-attained
\chapters@set{20261125T009000}{19}% kronecker-limit-and-chowla-selberg
\chapters@set{20261125T002000}{20}% definite-theta-lifts
\chapters@set{20261125T015500}{21}% siegel-weil-in-the-smallest-case
\chapters@set{20261125T002100}{22}% definite-theta-lifts-and-norms
\chapters@set{20261125T016500}{23}% hecke-operators-and-theta-lifts
\chapters@set{20261125T016700}{24}% divisibility-along-the-chain
\chapters@set{20261125T004000}{25}% denominators-and-sha
\chapters@set{20261125T004100}{26}% tate-twists-and-sha
\chapters@set{20261126T020430}{27}% the-minus-part-by-gauss-sums
\chapters@set{20261125T014100}{28}% congruence-modules
\chapters@set{20261125T014500}{29}% eisenstein-congruences-and-selmer-classes
\chapters@set{20261125T014700}{30}% pseudocharacters-and-the-length-bound
\chapters@set{20261125T011000}{31}% galois-representations-black-box
\chapters@set{20261126T030300}{32}% the-data
\chapters@set{20261126T020210}{33}% bessel-poincare-series-and-their-gram-operators
\chapters@set{20261126T030400}{34}% the-seesaw-identity
\chapters@set{20261126T030500}{35}% divisibility-of-the-composite-lift
\chapters@set{20261126T030600}{36}% theorem-a-short-forms
\chapters@set{20261126T030700}{37}% the-determinant-of-a-chain
\chapters@set{20261126T030800}{38}% the-cost-of-one-lift
\chapters@set{20261126T030900}{39}% denominators-of-projections-to-lifts
\chapters@set{20261126T031000}{40}% twists-with-few-congruences
\chapters@set{20261126T031100}{41}% theorem-c-and-the-main-inequality
\providecommand\chapnumber[1]{%
  \@ifundefined{chapters@#1}%
    {\PackageError{chapters}{Sheet #1 is not listed in the index}%
      {List it in 20261125T012000--grh-workbook__guide.org\MessageBreak
       and run scripts/make-chapters.py.}}%
    {\edef\thechap{\@nameuse{chapters@#1}}}}
\makeatother
\chapnumber{<sheet ID>}
\makeatletter
\providecommand\chapters@set[2]{\@namedef{chapters@#1}{#2}}
\chapters@set{20261125T001000}{1}% zeros-to-ramp-sums
\chapters@set{20261125T001100}{2}% inert-primes-and-estimate-b
\chapters@set{20261125T005000}{3}% galois-fourier-and-averaging
\chapters@set{20261125T005500}{4}% columns-from-a-sphere
\chapters@set{20261125T000000}{5}% slope-method
\chapters@set{20261126T030100}{6}% a-twist-for-each-prime
\chapters@set{20261125T003000}{7}% the-chain-and-the-square-root
\chapters@set{20261125T010000}{8}% rank-lemma-and-zero-estimate
\chapters@set{20261125T010500}{9}% poincare-series-and-near-orthogonality
\chapters@set{20261124T223000}{10}% elliptic-curves-for-the-proof
\chapters@set{20261125T009500}{11}% fock-model-and-theta-intertwiners
\chapters@set{20261124T223100}{12}% complex-multiplication-and-inert-reduction
\chapters@set{20261125T008000}{13}% one-cm-determinant
\chapters@set{20261124T214500}{14}% supersingular-theta-divisibility
\chapters@set{20261124T214600}{15}% the-isotropic-plane
\chapters@set{20261125T013700}{16}% pairings-on-a-formal-group-of-height-two
\chapters@set{20261125T013500}{17}% hecke-nondegeneracy-on-the-tree
\chapters@set{20261126T030200}{18}% the-bound-is-attained
\chapters@set{20261125T009000}{19}% kronecker-limit-and-chowla-selberg
\chapters@set{20261125T002000}{20}% definite-theta-lifts
\chapters@set{20261125T015500}{21}% siegel-weil-in-the-smallest-case
\chapters@set{20261125T002100}{22}% definite-theta-lifts-and-norms
\chapters@set{20261125T016500}{23}% hecke-operators-and-theta-lifts
\chapters@set{20261125T016700}{24}% divisibility-along-the-chain
\chapters@set{20261125T004000}{25}% denominators-and-sha
\chapters@set{20261125T004100}{26}% tate-twists-and-sha
\chapters@set{20261126T020430}{27}% the-minus-part-by-gauss-sums
\chapters@set{20261125T014100}{28}% congruence-modules
\chapters@set{20261125T014500}{29}% eisenstein-congruences-and-selmer-classes
\chapters@set{20261125T014700}{30}% pseudocharacters-and-the-length-bound
\chapters@set{20261125T011000}{31}% galois-representations-black-box
\chapters@set{20261126T030300}{32}% the-data
\chapters@set{20261126T020210}{33}% bessel-poincare-series-and-their-gram-operators
\chapters@set{20261126T030400}{34}% the-seesaw-identity
\chapters@set{20261126T030500}{35}% divisibility-of-the-composite-lift
\chapters@set{20261126T030600}{36}% theorem-a-short-forms
\chapters@set{20261126T030700}{37}% the-determinant-of-a-chain
\chapters@set{20261126T030800}{38}% the-cost-of-one-lift
\chapters@set{20261126T030900}{39}% denominators-of-projections-to-lifts
\chapters@set{20261126T031000}{40}% twists-with-few-congruences
\chapters@set{20261126T031100}{41}% theorem-c-and-the-main-inequality
\providecommand\chapnumber[1]{%
  \@ifundefined{chapters@#1}%
    {\PackageError{chapters}{Sheet #1 is not listed in the index}%
      {List it in 20261125T012000--grh-workbook__guide.org\MessageBreak
       and run scripts/make-chapters.py.}}%
    {\edef\thechap{\@nameuse{chapters@#1}}}}
\makeatother
\chapnumber{<sheet ID>}
\makeatletter
\providecommand\chapters@set[2]{\@namedef{chapters@#1}{#2}}
\chapters@set{20261125T001000}{1}% zeros-to-ramp-sums
\chapters@set{20261125T001100}{2}% inert-primes-and-estimate-b
\chapters@set{20261125T005000}{3}% galois-fourier-and-averaging
\chapters@set{20261125T005500}{4}% columns-from-a-sphere
\chapters@set{20261125T000000}{5}% slope-method
\chapters@set{20261126T030100}{6}% a-twist-for-each-prime
\chapters@set{20261125T003000}{7}% the-chain-and-the-square-root
\chapters@set{20261125T010000}{8}% rank-lemma-and-zero-estimate
\chapters@set{20261125T010500}{9}% poincare-series-and-near-orthogonality
\chapters@set{20261124T223000}{10}% elliptic-curves-for-the-proof
\chapters@set{20261125T009500}{11}% fock-model-and-theta-intertwiners
\chapters@set{20261124T223100}{12}% complex-multiplication-and-inert-reduction
\chapters@set{20261125T008000}{13}% one-cm-determinant
\chapters@set{20261124T214500}{14}% supersingular-theta-divisibility
\chapters@set{20261124T214600}{15}% the-isotropic-plane
\chapters@set{20261125T013700}{16}% pairings-on-a-formal-group-of-height-two
\chapters@set{20261125T013500}{17}% hecke-nondegeneracy-on-the-tree
\chapters@set{20261126T030200}{18}% the-bound-is-attained
\chapters@set{20261125T009000}{19}% kronecker-limit-and-chowla-selberg
\chapters@set{20261125T002000}{20}% definite-theta-lifts
\chapters@set{20261125T015500}{21}% siegel-weil-in-the-smallest-case
\chapters@set{20261125T002100}{22}% definite-theta-lifts-and-norms
\chapters@set{20261125T016500}{23}% hecke-operators-and-theta-lifts
\chapters@set{20261125T016700}{24}% divisibility-along-the-chain
\chapters@set{20261125T004000}{25}% denominators-and-sha
\chapters@set{20261125T004100}{26}% tate-twists-and-sha
\chapters@set{20261126T020430}{27}% the-minus-part-by-gauss-sums
\chapters@set{20261125T014100}{28}% congruence-modules
\chapters@set{20261125T014500}{29}% eisenstein-congruences-and-selmer-classes
\chapters@set{20261125T014700}{30}% pseudocharacters-and-the-length-bound
\chapters@set{20261125T011000}{31}% galois-representations-black-box
\chapters@set{20261126T030300}{32}% the-data
\chapters@set{20261126T020210}{33}% bessel-poincare-series-and-their-gram-operators
\chapters@set{20261126T030400}{34}% the-seesaw-identity
\chapters@set{20261126T030500}{35}% divisibility-of-the-composite-lift
\chapters@set{20261126T030600}{36}% theorem-a-short-forms
\chapters@set{20261126T030700}{37}% the-determinant-of-a-chain
\chapters@set{20261126T030800}{38}% the-cost-of-one-lift
\chapters@set{20261126T030900}{39}% denominators-of-projections-to-lifts
\chapters@set{20261126T031000}{40}% twists-with-few-congruences
\chapters@set{20261126T031100}{41}% theorem-c-and-the-main-inequality
\providecommand\chapnumber[1]{%
  \@ifundefined{chapters@#1}%
    {\PackageError{chapters}{Sheet #1 is not listed in the index}%
      {List it in 20261125T012000--grh-workbook__guide.org\MessageBreak
       and run scripts/make-chapters.py.}}%
    {\edef\thechap{\@nameuse{chapters@#1}}}}
\makeatother
\chapnumber{20261125T002000}
\renewcommand{\thesection}{\thechap.\arabic{section}}

\theoremstyle{definition}
\newtheorem{problem}{Problem}[section]
\theoremstyle{plain}
\newtheorem{theorem}{Theorem}[section]

\newcommand{\Z}{\mathbb{Z}}
\newcommand{\Q}{\mathbb{Q}}
\newcommand{\R}{\mathbb{R}}
\newcommand{\C}{\mathbb{C}}
\newcommand{\ii}{\mathrm{i}}
\newcommand{\Oc}{\mathcal{O}}
\newcommand{\Sc}{\mathbb{S}}
\newcommand{\Ac}{\mathcal{A}}
\newcommand{\Cl}{\mathrm{Cl}}
\newcommand{\Sym}{\operatorname{Sym}}
\newcommand{\tr}{\operatorname{tr}}

\setlength{\emergencystretch}{2em}

\title{Masses and the count of forms}
\author{}
\date{November 25, 2026}

\begin{document}
\maketitle
\refstepcounter{chap}\label{chap}

\paragraph{Prerequisites.}
Hermitian forms over the Gaussian integers, at the level of definitions.  The classification of positive definite unimodular \(\Z\)-lattices of rank at most \(8\) \cite[Ch.~16]{CoS}.  Schur functors and the Weyl dimension formula for \(\mathrm{GL}_n(\C)\) \cite[Lecture~6, Theorem~6.3]{FH}.  The values of Dirichlet \(L\)-functions at negative integers \cite[Theorem~4.2]{Wa}, and the functional equations of \(\zeta(s)\) and \(L(s,\chi)\) \cite[Ch.~8, 9]{Dav}.  Stirling's formula.  No earlier chapter is used.

In Alp\"oge's proof the maps of the chain of Chapter~\ref{ch-chap} are theta lifts between spaces of automorphic forms on definite unitary groups.  We take definite unitary groups over \(\Z[\ii]\) as a model.  A definite unitary group has finitely many lattice classes in a genus, and its automorphic forms are vectors attached to these classes.  How many forms are there?  The answer is the product of a mass and a Weyl dimension, and both factors carry \(\prod_{j<n} j!\), which cancels.  The theta lifts between these spaces are the subject of Chapter~\ref{tn-chap}, \emph{Definite theta lifts and their norms}; Chapter~\ref{sw-chap} first evaluates the genus averages they need.

\section{Hermitian lattices over \texorpdfstring{\(\Z[\ii]\)}{Z[i]}}\label{sec:lattices}

Let \(\Oc := \Z[\ii]\) and \(K := \Q(\ii)\).  A \emph{Hermitian lattice} of rank \(n\) is a free \(\Oc\)-module \(L\) of rank \(n\) with a positive definite Hermitian form \(h \colon L \times L \to \Oc\), linear in the first variable.  Its dual is \(L^\sharp := \{y \in L \otimes_\Oc K : h(x,y) \in \Oc \text{ for all } x \in L\}\), and \(L\) is \emph{unimodular} if \(L^\sharp = L\).  The basic example is \(I_n := \Oc^n\) with \(h(x,y) = \sum_k x_k \bar y_k\).  Let \(\Gamma_L\) be the group of \(\Oc\)-linear automorphisms of \(L\) preserving \(h\); it is finite, since it permutes the finitely many vectors of bounded length that span \(L\).

Two Hermitian lattices are in the same \emph{genus} if they become isometric over \(\Oc \otimes \Z_p\) for every prime \(p\).  Let \(\Cl_n\) be the set of isometry classes in the genus of \(I_n\).  By Problem~\ref{pr:real} below, a class gives a unimodular \(\Z\)-lattice of rank \(2n\) together with an automorphism of it; there are finitely many such lattices up to isometry \cite[Ch.~9]{Ca}, each with a finite automorphism group, so \(\Cl_n\) is finite.  It is the Shimura variety, of dimension zero, of the definite unitary group of \(I_n\); we call it the \emph{class set}.  Its \emph{mass} is
\[
  M_n := \sum_{L \in \Cl_n} \frac{1}{\lvert \Gamma_L \rvert}, \qquad M_0 := 1 .
\]
What do the class sets look like in small rank?

\begin{problem}\label{pr:real}
Let \(L\) be a Hermitian lattice and \(B(x,y) := \operatorname{Re} h(x,y)\).  Then \(B\) is a positive definite symmetric \(\Z\)-bilinear form with \(B(x,x) = h(x,x)\), multiplication by \(\ii\) preserves \(B\), and \(h(x,y) = B(x,y) + \ii B(x, \ii y)\).  The dual of \(L\) for \(B\) is \(L^\sharp\); in particular \(L\) is unimodular exactly when \((L, B)\) is.
\end{problem}

\begin{problem}\label{pr:aut}
\(\Gamma_{I_n}\) is the group of monomial matrices whose nonzero entries lie in \(\mu_4 = \{\pm 1, \pm \ii\}\).  Its order is \(4^n n!\).
\end{problem}

\begin{problem}\label{pr:odd}
Every lattice in the genus of \(I_n\) contains a vector \(x\) with \(h(x,x)\) odd.
\end{problem}

\begin{problem}\label{pr:class-one}
If \(n \le 4\), then \(\Cl_n = \{I_n\}\).
\end{problem}

Combining Problems~\ref{pr:aut} and~\ref{pr:class-one} gives \(M_n = 1/(4^n n!)\) for \(n \le 4\).  In rank \(5\) a second class appears: \(E \perp I_1\), where \(E\) is \(E_8\) with a Hermitian structure over \(\Oc\), in which \(h(x,x)\) is always even, and \(\lvert \Gamma_E \rvert = 46080\) is the order of the reflection group \(G_{31}\) of Shephard and Todd \cite{ShT}.  There are no others (Problem~\ref{pr:rank-five}).


\section{The mass of a genus}\label{sec:mass}

How does \(M_n\) grow?  Let \(\chi\) be the non-trivial character modulo \(4\), and let \(\chi^j\) be \(\chi\) for odd \(j\) and the trivial character for even \(j\), so that \(L(s, \chi^j) = \zeta(s)\) for even \(j\).  For the genus of \(I_n\), the mass formula reads
\begin{equation}\label{eq:mass}
  M_n = \delta_n \prod_{j=1}^{n} \frac{\lvert L(1-j, \chi^j) \rvert}{2},
\end{equation}
where \(\delta_n\) is a positive rational number, a local factor at the prime \(2\), which is ramified in \(K\).  This is the Hermitian analogue of Siegel's mass formula for quadratic forms \cite{Si}, with the local factors computed by Gan and Yu \cite{GY} and, at the ramified prime \(2\), by Cho \cite{Cho}; made explicit in \cite[Corollary~4.3]{JZ}, it gives \(\delta_n = 1\) for odd \(n\) and \(\delta_n = 2^n - 1\) for even \(n \ge 2\).  We set \(\delta_0 := 1\), so that \eqref{eq:mass} also holds for \(n = 0\), where \(M_0 = 1\).

\begin{problem}\label{pr:Lvalues}
For \(j = 1, 2, 3, 4, 5\), the numbers \(\lvert L(1-j, \chi^j) \rvert / 2\) are \(1/4\), \(1/24\), \(1/4\), \(1/240\), \(5/4\).
\end{problem}

\begin{problem}\label{pr:delta}
\(\delta_1 = \delta_3 = 1\), \(\delta_2 = 3\) and \(\delta_4 = 15\).
\end{problem}

In rank \(5\), with \(\delta_5 = 1\), the formula gives \(M_5 = 1/73728\).  Two unimodular lattices of the same rank over \(\Oc\), each containing a vector \(x\) with \(h(x,x)\) odd, lie in the same genus \cite[Corollary~2.2]{JZ}; thus \(E \perp I_1 \in \Cl_5\).

\begin{problem}\label{pr:rank-five}
The only vectors \(x \in E \perp I_1\) with \(h(x,x) = 1\) are the four multiples \(u e\), \(u \in \mu_4\), of a generator \(e\) of \(I_1\); hence \(\Gamma_{E \perp I_1} = \Gamma_E \times \mu_4\).  Moreover \(\Cl_5 = \{I_5, E \perp I_1\}\).
\end{problem}

The functional equation turns the values at \(1 - j\) into values at \(j\), and in doing so brings out a factorial.

\begin{problem}\label{pr:FE}
For \(j \ge 1\),
\[
  \frac{\lvert L(1-j, \chi^j) \rvert}{2} = (j-1)!\, \pi^{-j} \beta_j,
  \qquad
  \beta_j :=
  \begin{cases}
    2^{j-1} L(j, \chi) & j \text{ odd},\\
    2^{-j} \zeta(j) & j \text{ even}.
  \end{cases}
\]
\end{problem}

For \(r \ge 1\) define \(\kappa_r > 0\) by
\begin{equation}\label{eq:kappa}
  \frac{M_r}{M_{r-1}} = (r-1)!\, \pi^{-r} \kappa_r .
\end{equation}
By \eqref{eq:mass} and Problem~\ref{pr:FE}, \(\kappa_r = (\delta_r/\delta_{r-1})\beta_r\).  Multiplying \eqref{eq:kappa} over \(r \le n\),
\begin{equation}\label{eq:mass-product}
  M_n = \pi^{-n(n+1)/2} \prod_{j=0}^{n-1} j! \, \prod_{j=1}^{n} \kappa_j .
\end{equation}
So the mass is a product of factorials, up to a factor whose logarithm is linear in \(n\) at each step.

\begin{problem}\label{pr:kappa}
For \(r \le 4\), \(\kappa_r = \pi^r/(4 \cdot r!)\); thus \(\kappa_1, \dots, \kappa_4\) are \(\pi/4\), \(\pi^2/8\), \(\pi^3/24\), \(\pi^4/96\).  With \(\delta_5 = 1\), \(\kappa_5 = \pi^5/288\).
\end{problem}

\begin{problem}\label{pr:kappa-bound}
For every \(r \ge 1\), \(3/4 \le \kappa_r \le 2\); in particular \(\log \kappa_r = O(1)\).
\end{problem}

  By \eqref{eq:mass-product}, \(\log M_n\) grows like \(\tfrac12 n^2 \log n\); since every \(\Gamma_L\) contains \(\mu_4\), the class set has at least \(4M_n\) elements, more than \(8 \cdot 10^5\) for \(n = 15\).

\section{Forms and the Weyl dimension formula}\label{sec:forms}

On a definite group a form is a function on the class set, with values in a representation.  How many are there?

For \(L \in \Cl_n\) let \(L_\C := L \otimes_\Oc \C\), using the inclusion \(\Oc \subset \C\).  The form \(h\) extends to a Hermitian inner product on \(L_\C\), and \(\Gamma_L\) is a finite subgroup of its unitary group.  A \emph{weight} is a sequence of integers \(\lambda = (\lambda_1 \ge \dots \ge \lambda_n \ge 0)\); write \(\lvert \lambda \rvert := \sum_k \lambda_k\).  For a finite-dimensional Hermitian space \(A\), let \(\Sc_\lambda(A) \subset A^{\otimes \lvert \lambda \rvert}\) be the image of the Young symmetrizer \cite[\S6.1]{FH}, with the inner product induced from \(A^{\otimes \lvert\lambda\rvert}\).  A linear map \(x \colon A \to B\) induces \(\Sc_\lambda(x) \colon \Sc_\lambda(A) \to \Sc_\lambda(B)\), the restriction of \(x^{\otimes \lvert\lambda\rvert}\).  Then \(\Sc_\lambda(xy) = \Sc_\lambda(x)\Sc_\lambda(y)\) and \(\Sc_\lambda(x)^* = \Sc_\lambda(x^*)\), the latter because \((x^*)^{\otimes k}\) commutes with the symmetric group.  If \(\lambda\) has at most \(\dim A\) nonzero entries, \(\Sc_\lambda(A)\) is the irreducible representation of \(\mathrm{GL}(A)\) of highest weight \(\lambda\); otherwise it is zero \cite[Theorem~6.3]{FH}.  For \(\lambda = (a)\) it is \(\Sym^a(A)\).  There, for \(u \in A\) we write \(u^a := u \otimes \dots \otimes u\); for a basis \(e_1, \dots, e_n\) of \(A\) and \(\lvert\nu\rvert = a\), the vectors \(e^\nu\) are defined by \((\sum_k t_k e_k)^a = \sum_\nu (a!/\nu!)\, t^\nu e^\nu\), and \(e^\nu e^\mu := e^{\nu+\mu}\) is the product of the symmetric algebra.

Let \(V_\lambda(L) := \Sc_\lambda(L_\C)\).  Weyl's dimension formula reads, with \(z_k := \lambda_k + n - k\),
\begin{equation}\label{eq:weyl}
  \dim V_\lambda(L) = D_n(z) := \frac{V_n(z)}{\prod_{j=0}^{n-1} j!},
  \qquad
  V_n(z) := \prod_{1 \le k < l \le n} (z_k - z_l)
\end{equation}
\cite[Theorem~6.3]{FH}.  Fix one lattice in each class of \(\Cl_n\).  A \emph{form of weight \(\lambda\)} assigns to each \(L \in \Cl_n\) a vector \(\varphi(L) \in V_\lambda(L)\) fixed by \(\Gamma_L\).  The forms make up a space \(\Ac_\lambda(n)\), with the norm
\[
  \lVert \varphi \rVert^2 := \sum_{L \in \Cl_n} \frac{1}{\lvert\Gamma_L\rvert}\, \lVert \varphi(L) \rVert^2 .
\]
These are the algebraic modular forms of Gross \cite{Gr} on the unitary group of \(I_n\), of level one.

\begin{problem}\label{pr:weyl}
\(\dim \Sym^a(\C^n) = \binom{a+n-1}{n-1}\), and \(\dim \Sc_{(3,2,0)}(\C^3) = 15\).  The denominator in \eqref{eq:weyl} is \(\prod_{j=0}^{n-1} j! = V_n(n-1, \dots, 1, 0)\).
\end{problem}

\begin{problem}\label{pr:sym-norm}
In \(\Sym^a(A)\), \(\langle u^a, v^a \rangle = \langle u, v \rangle^a\).  If \(e_1, \dots, e_n\) is orthonormal, the \(e^\nu\) are orthogonal and \(\lVert e^\nu \rVert^2 = \nu!/a!\).
\end{problem}

\begin{problem}\label{pr:forms-small}
\(\Ac_\lambda(n) = 0\) unless \(4 \mid \lvert\lambda\rvert\).  The space \(\Ac_{(a)}(1)\) is one-dimensional when \(4 \mid a\).
\end{problem}

\begin{problem}\label{pr:dim-forms}
\(\dim \Ac_\lambda(n) = \sum_{L} \lvert\Gamma_L\rvert^{-1} \sum_{\gamma \in \Gamma_L} \tr(\gamma \mid V_\lambda(L))\).  When \(4 \mid \lvert\lambda\rvert\), the scalars \(\gamma \in \mu_4\) contribute \(4 M_n \dim V_\lambda\) to this sum.
\end{problem}

\begin{problem}\label{pr:forms-2}
The \(\Gamma_{I_2}\)-invariants in the symmetric algebra of \(\C^2\) form the polynomial algebra on \(e_1^4 + e_2^4\) and \(e_1^4 e_2^4\).  Hence, for \(4 \mid a\), \(\dim \Ac_{(a)}(2) = \lfloor a/8 \rfloor + 1\), which is \((1 + o(1)) \cdot 4 M_2 \dim V_{(a)}\) as \(a \to \infty\).
\end{problem}

Does \(\lvert\lambda\rvert \to \infty\) suffice?  Compare two weights of the same size.

\begin{problem}\label{pr:two-weights}
Let \(j \ge 1\), \(\lambda := (8j, 0)\) and \(\lambda' := (4j, 4j)\).  Then \(\dim \Ac_\lambda(2) = j + 1\) and \(\dim \Ac_{\lambda'}(2) = 1\), while the scalar contributions \(4 M_2 \dim V\) of Problem~\ref{pr:dim-forms} are \(j + \frac18\) and \(\frac18\).  In both cases the non-scalar elements of \(\Gamma_{I_2}\) contribute exactly \(\frac78\).
\end{problem}

The non-scalar elements contribute the same amount to both weights.  Beside \(\dim V_\lambda = 8j + 1\) it is negligible; for \(\lambda'\), where \(V_{\lambda'} = \det^{4j}\) is a line on which all of \(\Gamma_{I_2}\) acts trivially, it is seven times the scalar term.  What must grow is the gap \(\lambda_1 - \lambda_2\), not \(\lvert\lambda\rvert\).

The same holds in general.  When \(4 \mid \lvert\lambda\rvert\) and all the gaps \(\lambda_k - \lambda_{k+1}\) tend to infinity, the Weyl character formula shows that each non-scalar \(\gamma\) contributes \(o(\dim V_\lambda)\), so \(\dim \Ac_\lambda(n) \sim 4 M_n \dim V_\lambda\).  The count of forms is thus the product of the mass and the Weyl dimension.  By \eqref{eq:mass-product} and \eqref{eq:weyl} each factor contains \(\prod_{j<n} j!\), once in the numerator and once in the denominator; \S\ref{sec:count} returns to this.

\section{The count of forms}\label{sec:count}

When \(4 \mid \lvert\lambda\rvert\) and the gaps of \(\lambda\) are large, the space of forms of weight \(\lambda\) has dimension about \(4 M_n \dim V_\lambda\) (\S\ref{sec:forms}).  The mass contains the product \(\prod_{j<n} j!\), whose logarithm is about \(\tfrac12 n^2 \log n\), and the Weyl dimension contains its inverse.  The norm of a theta lift is controlled to within a term linear in the weight, so the count must be known to the same accuracy.  Why is it?

\begin{problem}\label{pr:stirling}
\(\sum_{j=0}^{n-1} \log j! = \tfrac12 n^2 \log n - \tfrac34 n^2 + O(n \log n)\).
\end{problem}

\begin{problem}\label{pr:count}
Let \(r_0 < r_1\), and let \(u \in \Z^{r_0}\) and \(w \in \Z^{r_1}\) be strictly decreasing.  Then
\[
  \frac{M_{r_1} D_{r_1}(w)}{M_{r_0} D_{r_0}(u)}
  = \frac{V_{r_1}(w)}{V_{r_0}(u)}\; \pi^{-\sum_{r = r_0+1}^{r_1} r} \prod_{r = r_0+1}^{r_1} \kappa_r .
\]
\end{problem}

The extra \((r-1)!\) in the mass ratio \eqref{eq:kappa} cancels the extra \((r-1)!\) in the denominator of \eqref{eq:weyl}.  With \(r_0 = 0\), where \(M_0 = D_0 = V_0 = 1\) (empty products), Problem~\ref{pr:count} reads as follows.

\begin{theorem}[The count identity]\label{thm:count}
For every \(r \ge 1\) and every strictly decreasing \(z \in \Z^r\),
\[
  M_r D_r(z) = V_r(z)\, \pi^{-r(r+1)/2} \prod_{j=1}^{r} \kappa_j .
\]
\end{theorem}

No factorial survives: what remains is a product of differences of weights, a power of \(\pi\) and the factors \(\kappa_j\).

The simplest instance adds one rank, with one large new weight.

\begin{problem}\label{pr:one-row}
Let \(u \in \Z^{r_0}\) be strictly decreasing, and \(w = (w_1, u_1, \dots, u_{r_0})\) with \(w_1 > u_1\).  Then
\[
  \log \frac{M_{r_0+1} D_{r_0+1}(w)}{M_{r_0} D_{r_0}(u)}
  = \sum_{i=1}^{r_0} \bigl( \log(w_1 - u_i) - \log \pi \bigr) - \log\pi + \log \kappa_{r_0+1} .
\]
\end{problem}

Thus each new rank contributes, for every existing index \(i\), the quantity \(\log(w_1 - u_i) - \log\pi\), and since \(\log\kappa_r = O(1)\) (Problem~\ref{pr:kappa-bound}) the remaining term is of lower order when \(w_1\) is large.

\paragraph{In the Lean formalization.}
In the rewritten Lean proof \cite{GL} the count is used in the following form, in \texttt{New/Spanning/SecondMoment/SpanningTheorem.lean} of the readable edition.  On a definite space of rank \(r\) over the CM field of the proof, with a level small enough that no class has a nontrivial automorphism, the forms of the principal level have dimension \(\#\bar G \cdot \mathrm{mass}_r\) times the product over the infinite places of the formal degrees of the weights, where \(\bar G\) is a finite group whose order is the index of this level in the maximal one (\texttt{finrank\_principalForms\_spaceW}).  The formal degrees are Weyl dimensions, as in Section~\ref{sec:forms}.  Hence the dimension of the space of target forms of the proof is at most \(\#\bar G\) times this product times \(\mathrm{mass}_r\) (\texttt{targetDim\_le\_card\_mul\_besselDegree\_mul\_mass}).

\section*{Hints}

\subsection*{\S\ref{sec:lattices}}
\begin{description}
\item[\ref{pr:real}] Compute \(B(\ii x, \ii y)\) and \(B(x, \ii y)\).  For the dual, replace \(x\) by \(\ii x\) in the condition \(B(x,y) \in \Z\).
\item[\ref{pr:aut}] An automorphism permutes the vectors with \(h(x,x) = 1\).  Which vectors of \(I_n\) are these?
\item[\ref{pr:odd}] Over \(\Z_2\) the lattice is isometric to \(I_n \otimes \Z_2\).  Approximate a vector of odd length \(2\)-adically by a vector of \(L\).
\item[\ref{pr:class-one}] Use Problems~\ref{pr:real} and~\ref{pr:odd} and the classification of odd unimodular lattices of rank at most \(8\); then look at the vectors of length one.
\end{description}

\subsection*{\S\ref{sec:mass}}
\begin{description}
\item[\ref{pr:Lvalues}] \(B_{j,\chi} = 4^{j-1}\sum_{a=1}^{4} \chi(a) B_j(a/4)\), and \(\zeta(1-j) = -B_j/j\).
\item[\ref{pr:delta}] Compare \eqref{eq:mass} with \(M_n = 1/(4^n n!)\).
\item[\ref{pr:FE}] For \(\zeta\) use \(\zeta(1-s) = 2(2\pi)^{-s}\cos(\pi s/2)\Gamma(s)\zeta(s)\).  For the odd character \(\chi\) of conductor \(4\), whose Gauss sum is \(2\ii\), complete \(L(s,\chi)\) by \((4/\pi)^{(s+1)/2}\Gamma(\tfrac{s+1}{2})\), and use the reflection and duplication formulas for \(\Gamma\).  Check the case \(j = 1\): \(L(1, \chi) = \pi/4\).
\item[\ref{pr:rank-five}] \(E\) is even.  For the second part, compare with the mass \(M_5\).
\item[\ref{pr:kappa}] Use \(M_r/M_{r-1} = 1/(4r)\) for \(r \le 4\), and \(M_5 = 1/73728\).
\item[\ref{pr:kappa-bound}] Use \(\kappa_r = (\delta_r/\delta_{r-1})\beta_r\) and the values of \(\delta_n\), treating odd and even \(r\) separately; \(1 - 3^{-r} \le L(r,\chi) \le 1\).  For \(r = 1\) use Problem~\ref{pr:kappa}.
\end{description}

\subsection*{\S\ref{sec:forms}}
\begin{description}
\item[\ref{pr:weyl}] For \(\lambda = (a)\), \(z = (a + n - 1, n - 2, \dots, 0)\).
\item[\ref{pr:sym-norm}] Expand \(\langle (\sum t_k e_k)^a, (\sum s_k e_k)^a \rangle\) in two ways.
\item[\ref{pr:forms-small}] The scalar \(\ii \in \Gamma_L\) acts on \(V_\lambda(L)\) by \(\ii^{\lvert\lambda\rvert}\).
\item[\ref{pr:dim-forms}] The dimension of the invariants of a finite group is the average of the trace.
\item[\ref{pr:forms-2}] An invariant must be invariant under \(e_k \mapsto \ii e_k\) and under the swap.  Count the pairs \((i,j)\) with \(4i + 8j = a\).
\item[\ref{pr:two-weights}] For \(\lambda\) use Problem~\ref{pr:forms-2}.  For \(\lambda'\), \(\Sc_{(k,k)}(\C^2)\) is the \(k\)-th power of the determinant, and every element of \(\Gamma_{I_2}\) has determinant in \(\mu_4\).  Use \(M_2 = 1/32\); the non-scalar part is the dimension less the scalar contribution.
\end{description}
\subsection*{\S\ref{sec:count}}
\begin{description}
\item[\ref{pr:stirling}] \(\log j! = j\log j - j + O(\log(j+1))\); compare the sum with an integral.
\item[\ref{pr:count}] Write the left side as a telescoping product, and use \eqref{eq:kappa} and \eqref{eq:weyl}.
\item[\ref{pr:one-row}] \(V_{r_0+1}(w) = V_{r_0}(u) \prod_i (w_1 - u_i)\).
\end{description}

\begin{thebibliography}{CoS}
\bibitem[Ca]{Ca} J.~W.~S.~Cassels, \emph{Rational Quadratic Forms}, Academic Press, 1978.
\bibitem[Cho]{Cho} S.~Cho, \emph{Group schemes and local densities of ramified hermitian lattices in residue characteristic 2}, arXiv:1210.7894, 2012.
\bibitem[CoS]{CoS} J.~H.~Conway and N.~J.~A.~Sloane, \emph{Sphere Packings, Lattices and Groups}, 3rd ed., Grundlehren Math.\ Wiss.\ 290, Springer, 1999.
\bibitem[Dav]{Dav} H.~Davenport, \emph{Multiplicative Number Theory}, 3rd ed., Grad.\ Texts in Math.\ 74, Springer, 2000.
\bibitem[FH]{FH} W.~Fulton and J.~Harris, \emph{Representation Theory: A First Course}, Grad.\ Texts in Math.\ 129, Springer, 1991.
\bibitem[GY]{GY} W.~T.~Gan and J.-K.~Yu, \emph{Group schemes and local densities}, Duke Math.\ J.\ \textbf{105} (2000), 497--524.
\bibitem[GL]{GL} L.~Sernau, \emph{GRH for Dirichlet \(L\)-functions, in Lean~4}, a shorter rewriting of the Lean proof, text by Claude, \url{https://github.com/lsernau/grh-lean}, commit \texttt{72fead1}, readable edition \texttt{grh/}.
\bibitem[Gr]{Gr} B.~H.~Gross, \emph{Algebraic modular forms}, Israel J.\ Math.\ \textbf{113} (1999), 61--93.
\bibitem[ShT]{ShT} G.~C.~Shephard and J.~A.~Todd, \emph{Finite unitary reflection groups}, Canad.\ J.\ Math.\ \textbf{6} (1954), 274--304.
\bibitem[JZ]{JZ} M.~J\"urgens and M.~C.~Zimmermann, \emph{Steinitz classes and partial genera of unimodular lattices over imaginary-quadratic fields}, arXiv:1512.03297, 2015.
\bibitem[Si]{Si} C.~L.~Siegel, \emph{\"Uber die analytische Theorie der quadratischen Formen}, Ann.\ of Math.\ \textbf{36} (1935), 527--606.
\bibitem[Wa]{Wa} L.~C.~Washington, \emph{Introduction to Cyclotomic Fields}, 2nd ed., Grad.\ Texts in Math.\ 83, Springer, 1997.
\end{thebibliography}

\end{document}
